\begin{prerequis}
Lire les coordonnées d'un point; reconnaître parallélogramme, translation,
symétrie centrale et milieu d'un segment.
\end{prerequis}
\begin{automatismes}
Dans un repère, avec $A(1;2)$ et $B(4;-1)$, calculer les coordonnées de
$\vv{AB}$ et du milieu de $[AB]$; placer le point $C$ tel que
$\vv{AC}=2\vv{AB}$.
\end{automatismes}
\begin{activite}[Activité d'entrée -- Décrire un déplacement]
Sur un plan quadrillé, plusieurs chemins conduisent du point $A$ au point $B$.
Coder chacun par ses déplacements horizontal et vertical. Identifier ce qui
reste invariant malgré le chemin choisi, puis proposer une règle d'addition de
deux déplacements successifs.
\end{activite}

